\anexo{Roteiro da Experiência}
\label{an:roteiro}

\textbf{``Segurança no Canteiro de Obras: Prevenção de Quedas e Uso de EPIs''}

\textit{Fase 1 do Sistema Modular de Treinamento em Realidade Virtual}

\section*{1. Introdução (Cena Inicial)}

\begin{enumerate}
    \item O usuário é recebido em um canteiro de obras virtual por um instrutor virtual (engenheiro de segurança do trabalho).
    \item Breve apresentação sobre a importância das normas de segurança NR-18 e NR-35, contextualizando os riscos de trabalho em altura.
    \item Instrução sobre como interagir com o ambiente virtual: movimentação, manipulação de objetos e sistema de \textit{feedback}.
\end{enumerate}

\section*{2. Seleção de EPIs (Cena 1)}

\textbf{Cenário:} O usuário se encontra em uma área de apoio do canteiro, diante de uma bancada com diversos equipamentos, incluindo EPIs corretos e incorretos para a atividade de trabalho em altura.

\textbf{Objetivos de aprendizagem:}

\begin{itemize}
    \item Identificar os EPIs obrigatórios para trabalho em altura conforme NR-35
    \item Distinguir EPIs adequados de equipamentos inadequados ou com defeito
\end{itemize}

\textbf{Interatividade:}

\begin{enumerate}
    \item O usuário deve selecionar os EPIs corretos para a atividade que será realizada: cinto de segurança tipo paraquedista, capacete com jugular, luvas e botina de segurança.
    \item Caso selecione um EPI incorreto ou com defeito visível, o sistema fornece \textit{feedback} imediato explicando o problema.
    \item Após a seleção correta, o instrutor virtual confirma os itens escolhidos e o usuário avança para o próximo cenário já equipado.
\end{enumerate}

\textbf{Pontuação:} Acerto na seleção de cada EPI obrigatório; penalização por seleção de item inadequado.

\textbf{Mensagem educativa:} \textit{``Antes de iniciar qualquer trabalho em altura, verifique se todos os EPIs estão em bom estado e são adequados para a atividade. Nunca substitua um EPI certificado por improvisações.''}

\section*{3. Trabalho em Altura (Cena 2)}

\textbf{Cenário:} O usuário se encontra em um andaime de uma edificação em construção. O ambiente contém irregularidades deliberadas: ausência de guarda-corpo em um trecho, cinto de segurança desconectado do ponto de ancoragem e rodapé ausente em parte da plataforma.

\textbf{Objetivos de aprendizagem:}

\begin{itemize}
    \item Identificar riscos de queda em ambiente de trabalho em altura
    \item Reconhecer ausências e irregularidades em EPCs
    \item Demonstrar o procedimento correto de conexão do cinto ao ponto de ancoragem certificado
\end{itemize}

\textbf{Interatividade:}

\begin{enumerate}
    \item O usuário deve inspecionar o andaime e identificar cada irregularidade presente, interagindo com os elementos de risco.
    \item Para cada risco identificado, o usuário deve executar a correção correspondente: instalar o guarda-corpo, conectar o cinto ao ponto de ancoragem e sinalizar a ausência do rodapé.
    \item Caso o usuário tente avançar sem corrigir as irregularidades, o sistema simula a consequência do risco: animação de queda ou queda de objeto, seguida de \textit{feedback} educativo explicando o que deveria ter sido feito.
    \item Após a correção de todas as irregularidades, o instrutor virtual apresenta um resumo das boas práticas demonstradas.
\end{enumerate}

\textbf{Pontuação:} Pontos por identificação e correção de cada irregularidade; penalização por tentativa de avanço sem correção.

\textbf{Mensagem educativa:} \textit{``Em trabalhos em altura, nunca inicie a atividade sem verificar a integridade dos EPCs e a correta fixação do cinto de segurança. Um único ponto de ancoragem não verificado pode custar uma vida.''}

\section*{4. Conclusão e Avaliação Final}

\begin{enumerate}
    \item O instrutor virtual apresenta um resumo das ações realizadas pelo usuário ao longo da experiência, destacando acertos e pontos de melhoria.
    \item O usuário recebe seu relatório de desempenho gerado automaticamente pelo \textit{SessionLogger}, contendo pontuação total, tarefas concluídas e tempo de execução.
    \item O instrutor reforça as principais diretrizes das NR-18 e NR-35 aplicadas nos cenários.
    \item O usuário recebe um certificado virtual de conclusão da Fase 1 do treinamento.
\end{enumerate}

\section*{5. Nota sobre Expansão Futura}

Este roteiro corresponde à Fase 1 do sistema modular. Fases futuras preveem a incorporação de cenários adicionais como queda de objetos, escavações e contenções, e inspeção geral de canteiro, a serem desenvolvidos e validados em parceria com especialistas e entidades do setor da construção civil.
